\section{August 28, 2026}

\subsection{Examples of vector spaces}

\begin{example}[Coordinate spaces]
  Let $\bb{F}$ be a field. The set
  \[
    \bb{F}^n
    =\left\{
      \begin{bmatrix}
        x_1\\ \vdots\\ x_n
      \end{bmatrix}
      :x_1,\ldots,x_n\in\bb{F}
    \right\}
  \]
  is a vector space over $\bb{F}$ under componentwise addition and scalar
  multiplication:
  \[
    \begin{bmatrix}x_1\\ \vdots\\ x_n\end{bmatrix}
    +\begin{bmatrix}y_1\\ \vdots\\ y_n\end{bmatrix}
    =\begin{bmatrix}x_1+y_1\\ \vdots\\ x_n+y_n\end{bmatrix},
    \qquad
    \alpha\begin{bmatrix}x_1\\ \vdots\\ x_n\end{bmatrix}
    =\begin{bmatrix}\alpha x_1\\ \vdots\\ \alpha x_n\end{bmatrix}.
  \]
  In particular, $\bb{R}^n$ is a vector space over $\bb{R}$ and
  $\bb{C}^n$ is a vector space over $\bb{C}$.
\end{example}

\begin{example}[Matrix spaces]
  The set $M_{m\times n}(\bb{F})$ of $m\times n$ matrices with entries in
  $\bb{F}$ is a vector space over $\bb{F}$. Addition and scalar
  multiplication are defined entrywise:
  \[
    [a_{ij}]+[b_{ij}]=[a_{ij}+b_{ij}],
    \qquad
    \alpha[a_{ij}]=[\alpha a_{ij}].
  \]
\end{example}

\begin{example}[Polynomial spaces]
  Let
  \[
    \mathcal P_n(\bb{F})
    =\{a_0+a_1t+\cdots+a_nt^n:a_0,\ldots,a_n\in\bb{F}\}.
  \]
  This is the vector space over $\bb{F}$ of polynomials of degree at most
  $n$. Its operations are
  \begin{align*}
    &(a_0+a_1t+\cdots+a_nt^n)
      +(b_0+b_1t+\cdots+b_nt^n)\\
    &\qquad=(a_0+b_0)+(a_1+b_1)t+\cdots+(a_n+b_n)t^n,
  \end{align*}
  and
  \[
    \alpha(a_0+a_1t+\cdots+a_nt^n)
    =\alpha a_0+\alpha a_1t+\cdots+\alpha a_nt^n.
  \]
\end{example}

\subsection{Linear combinations and spanning sets}

\begin{definition}[Linear combination and span]
  Let $V$ be a vector space over $\bb{F}$ and let
  $v_1,\ldots,v_p\in V$. A \emph{linear combination} of these vectors is
  a vector of the form
  \[
    \alpha_1v_1+\cdots+\alpha_pv_p,
    \qquad
    \alpha_1,\ldots,\alpha_p\in\bb{F}.
  \]
  The set of all their linear combinations is called their \emph{span}:
  \[
    \operatorname{span}\{v_1,\ldots,v_p\}
    =\{\alpha_1v_1+\cdots+\alpha_pv_p:
      \alpha_1,\ldots,\alpha_p\in\bb{F}\}.
  \]
  A linear combination is \emph{trivial} if all of its coefficients are
  zero.
\end{definition}

In $\bb{R}^2$, let
\[
  e_1=\begin{bmatrix}1\\0\end{bmatrix},
  \qquad
  e_2=\begin{bmatrix}0\\1\end{bmatrix}.
\]
Then
\[
  \operatorname{span}\{e_1,e_2\}=\bb{R}^2,
\]
because
\[
  \begin{bmatrix}x\\y\end{bmatrix}=xe_1+ye_2.
\]
If
\[
  v=\begin{bmatrix}1\\1\end{bmatrix},
\]
then $\{e_1,e_2,v\}$ also spans $\bb{R}^2$, since
\[
  \begin{bmatrix}x\\y\end{bmatrix}=xe_1+ye_2+0v.
\]
The additional vector is unnecessary, which motivates the question of
redundancy in a spanning set.

\begin{definition}[Spanning set]
  A collection $\{v_1,\ldots,v_p\}\subseteq V$ is a \emph{spanning set}
  for $V$ if every $v\in V$ can be written as a linear combination
  \[
    v=\alpha_1v_1+\cdots+\alpha_pv_p.
  \]
  Equivalently,
  \[
    \operatorname{span}\{v_1,\ldots,v_p\}=V.
  \]
\end{definition}

\subsection{Linear independence}

\begin{definition}[Linear independence]
  A collection $\{v_1,\ldots,v_p\}\subseteq V$ is \emph{linearly
  independent} if
  \[
    \alpha_1v_1+\cdots+\alpha_pv_p=0
  \]
  implies
  \[
    \alpha_1=\cdots=\alpha_p=0.
  \]
  In other words, only the trivial linear combination of
  $v_1,\ldots,v_p$ equals zero. A collection that is not linearly
  independent is called \emph{linearly dependent}.
\end{definition}

\begin{example}
  Consider the vectors
  \[
    v_1=\begin{bmatrix}1\\1\\0\end{bmatrix},
    \qquad
    v_2=\begin{bmatrix}2\\0\\1\end{bmatrix},
    \qquad
    v_3=\begin{bmatrix}3\\1\\1\end{bmatrix}
  \]
  in $\bb{R}^3$. To test linear independence, solve
  \[
    \alpha v_1+\beta v_2+\gamma v_3=0.
  \]
  This gives the system
  \[
    \begin{cases}
      \alpha+2\beta+3\gamma=0,\\
      \alpha+\gamma=0,\\
      \beta+\gamma=0.
    \end{cases}
  \]
  Taking $\gamma=1$ gives $\alpha=\beta=-1$, so
  \[
    -v_1-v_2+v_3=0.
  \]
  Thus the vectors are linearly dependent; equivalently,
  $v_3=v_1+v_2$.
\end{example}

The following result precisely characterizes redundancy in a finite
collection of vectors.

\begin{proposition}\label{prop:dependence-redundancy}
  Let $v_1,\ldots,v_p\in V$. The collection
  $\{v_1,\ldots,v_p\}$ is linearly dependent if and only if one of its
  vectors is a linear combination of the others. Equivalently, there is
  an index $k$ and scalars $\beta_j\in\bb{F}$ such that
  \[
    v_k=\sum_{\substack{j=1\\j\neq k}}^p\beta_jv_j.
  \]
\end{proposition}

\begin{proof}
  Suppose first that $\{v_1,\ldots,v_p\}$ is linearly dependent. Then
  there are scalars $\alpha_1,\ldots,\alpha_p$, not all zero, such that
  \[
    \alpha_1v_1+\cdots+\alpha_pv_p=0.
  \]
  Choose $k$ such that $\alpha_k\neq0$. Since $\bb{F}$ is a field,
  $\alpha_k^{-1}$ exists, and we can solve for $v_k$:
  \[
    v_k
    =-\sum_{\substack{j=1\\j\neq k}}^p
      \alpha_k^{-1}\alpha_jv_j.
  \]
  Thus $v_k$ is a linear combination of the other vectors.

  Conversely, suppose that for some $k$,
  \[
    v_k=\sum_{\substack{j=1\\j\neq k}}^p\beta_jv_j.
  \]
  Moving all terms to one side gives
  \[
    v_k-\sum_{\substack{j=1\\j\neq k}}^p\beta_jv_j=0.
  \]
  This is a nontrivial linear combination of $v_1,\ldots,v_p$, because
  the coefficient of $v_k$ is $1\neq0$. Hence the collection is linearly
  dependent.
\end{proof}
